% Page 014 — page imprimée 10
% ROMAN (Jean), vers 1668 - 1715
% Mise en page en deux colonnes, fidèle à l'original.


\begin{center}
  {\fontsize{12}{13}\selectfont\bfseries ROMAN (Jean)\par}
  {\fontsize{12}{13}\selectfont\bfseries vers 1668 - 1715\par}
  \vspace{0.10em}
  \rule{3.8cm}{0.45pt}
\end{center}

\vspace{0.45em}

{\fontsize{9.65}{11.05}\selectfont
\begin{multicols}{2}

Jean Roman occupe l’une des premières places dans la saga des prédicants protestants des Cévennes et du Bas-Languedoc, entre 1685 et 1701, à laquelle le pasteur et historien Charles Bost a consacré en 1912 deux forts volumes qui restent à ce jour un monument indépassable. Actif surtout dans les hautes Cévennes, entre Aigoual et Mont Lozère, il n’était pourtant pas, contrairement à la quasi totalité de ses compagnons d’apostolat, originaire de la région : dauphinois, il était né à Vercheny, dans la vallée de la Drôme, et avait gagné le Refuge à la Révocation de l’Edit de Nantes (1685). Installé pendant deux ans à Lausanne, il fut très impressionné par la lecture de l’ouvrage de Josué de la Place, \textit{Dialogue d’un père avec son fils sur la question si l’on peut faire son salut en allant à la messe}, et finit par prendre la résolution de revenir en France porter la bonne Parole à ses coreligionnaires. Entré par le Dauphiné, il résolut de n’y pas rester afin d’éviter les pressions de ses parents qui chercheraient à le détourner de sa périlleuse carrière, et prit « le chemin des Cévennes, où il avais oui dire qu’on faisait des assemblées ». Déjà les Cévennes s’imposaient à l’attention du protestantisme européen; et c’est sur leurs crêtes et dans leurs vallées tourmentées que le dauphinois allait se mêler intimement au peuple calviniste entré dans l’ère de la persécution, et vivre la plus extraordinaire des aventures, à coups d’assemblées, de fuites, de blessures, de trahisons et d’arrestations. Il ne manquera même pas à ce héros de roman des amours nombreuses avec de jeunes cévenoles (vertueusement fustigées par l’abbé du Chaïla, son ennemi juré!).

Si le vallerauguais Jean Vivent et l’avocat nîmois Claude Brousson furent les plus grands et restent les plus connus de ces prédicants, Jean Roman fut peut-être le plus attachant, et en tous les cas le plus fidèle; deux fois échappé, dans des conditions rocambolesques, des griffes de la justice royale, il n’eut pas à connaître le martyre. Mais il eut le courage, avec le gévaudanais Olivier (1), de poursuivre son apostolat jusqu’à l’extrême limite de ses forces, avant de devoir gagner définitivement le Refuge.

Roman abrita longtemps son apostolat sous des allures de colporteur: aussi fut-il rapidement connu sous les surnoms de Marchandou (le petit marchand), Le petit paquetier ou Lou biquarel (le colporteur). Rentré en France à la fin de 1687, il avait traversé le Vivarais et était venu se fixer, vers la fin de 1688, dans une région qui lui fut toujours chère, entre Vialas et Vébron, autour des montagnes du Lozère et du Bougès (le territoire même de l’abbé du Chaïla). Instruit (il savait le latin), il commença à tenir des assemblées clandestines. Il fut surpris une première fois sur le Mont Lozère, dans la nuit du 18 décembre 1688, mais échappa aux coups de baïonnettes des soldats qui sondaient le foin du grenier où il s’était réfugié, ses habits sous le bras.

\section*{Il s’évade de Montvaillant}

Quelques mois plus tard, alors qu’il se cachait chez l’un de ses fidèles, après une assemblée sur la Can de l’Hospitalet, il fut livré par un apostat aux soldats de Vébron qui le traînèrent au fort d’Alès, en le frappant durement tout au long du chemin. Transféré à Saint-Jean du Gard sur l’ordre de l’intendant Basville, qui voulait l’interroger lui-même, le 3 octobre 1689, il fut attaché sur un cheval sous le ventre duquel on passa une pièce de bois qui lui maintenait les pieds, et enfermé au château de M. de Montvaillant. Refusant de trahir des noms, Roman fut menacé d’être pendu par le commandant militaire de la province, le comte de Broglie, mais lui reprocha, en guise de réponse, de « se faire compagnon de la Bête enivrée du sang des martyrs de Jésus ». Jeté dans un cachot, garrotté de cordes des pieds aux épaules, il eut la surprise d’entendre par un trou du mur la voix d’une jeune fille, Suzanne Guichard (née à Aulas en 1671, fille du pasteur Lévy Guichard, mort en 1681, et de Suzanne Atger, de l’hermet de Saint-Julien d’Arpaon, elle était entrée au service des de Montvaillant); elle réussit à trancher avec la pointe d’une baïonnette les cordes de ses mains. Le prédicant se défit de ses autres liens avec la même arme, et réussit à creuser la pierre dans laquelle entrait le verrou du cachot.

Il dut alors, dans un château garni de troupes et où logeait l’intendant en personne, s’enfuir à l’aide de draps, en passant devant les fenêtres de Basville; ce dernier, le lendemain, croyant à une complicité de M. de Montvaillant (un nouveau converti), le menaça des pires maux. Suzanne Guichard s’était enfuie, mais apprenant les soupçons qui pesaient sur ses maîtres, elle revint au château et y rétablit publiquement la vérité. Condamnée au fouet et à la prison perpétuelle, el-

\end{multicols}
}

